% =====================================================================
% QUINTA Y SEXTA PÁGINA: aplicación de la metodología con datos y análisis
% de los resultados correspondientes.
% =====================================================================

\section*{Aplicación de la metodología}

\subsection*{Datos}

Se aplicó la metodología del artículo a un conjunto de datos sintético de consumo y clima (datos diarios de 20 países, 2020 a 2024). Siguiendo los lineamientos del artículo, se trabajó en frecuencia mensual y se eligió el país más pequeño del conjunto: Países Bajos. El proceso de limpieza deja un conjunto depurado mediante los siguientes pasos:

\begin{itemize}
  \item Filtrado por país, con la fecha como índice temporal, verificación de continuidad diaria y de valores faltantes (no había).
  \item Agregación mensual: la variable objetivo es el pico de consumo mensual, es decir, el máximo diario de consumo en cada mes (análogo al pico de demanda del artículo); como exógenas se usan la temperatura y la humedad medias del mes.
  \item Winsorización mensual (percentiles 5 y 95) sobre las series agregadas, de modo que solo se recorten los meses extremos.
\end{itemize}

\noindent El resultado son 60 observaciones mensuales, de enero de 2020 a diciembre de 2024.

\subsection*{Metodología aplicada}

Se replicaron los pasos del artículo: test ADF de estacionariedad (diferenciando si $p>0{,}05$), selección de órdenes por AIC, validación walk-forward sobre 2024 (horizonte $h=1$, reestimando en cada paso) y métricas MAE, RMSE y MAPE.

\begin{itemize}
  \item ARIMA: los órdenes $(p,d,q)$ se seleccionaron automáticamente minimizando el AIC, con estacionalidad anual. El modelo elegido fue ARIMA$(0,1,1)(1,1,0)_{12}$, con AIC $510{,}7$. El test de Ljung-Box sobre los residuos arrojó $p=0{,}0007$, por lo que aún queda autocorrelación residual por capturar.
  \item VAR: modelo multivariado sobre el pico de consumo, la temperatura y la humedad (diferenciadas para garantizar estacionariedad conjunta). El rezago se eligió por AIC y resultó ser 6 (AIC $12{,}6$). El Ljung-Box sobre los residuos del pico dio $p=0{,}62$, es decir, sin autocorrelación residual.
\end{itemize}

\subsection*{Resultados}

\begin{table}[H]
\centering
\caption{Desempeño comparativo de ARIMA y VAR sobre el pico mensual (Países Bajos, validación walk-forward 2024).}
\label{tab:app}
\begin{tabular}{lccccc}
\toprule
Modelo & Órdenes / rezago & AIC & MAE & RMSE & MAPE (\%) \\
\midrule
ARIMA & $(0,1,1)(1,1,0)_{12}$ & 510{,}7 & 324{,}71 & 373{,}99 & 2{,}64 \\
VAR   & rezago 6 & 12{,}6 & 345{,}09 & 463{,}83 & 3{,}26 \\
\bottomrule
\end{tabular}
\end{table}

\begin{figure}[H]
\centering
\includegraphics[width=0.82\linewidth]{arima_forecast.png}
\caption{Pronóstico ARIMA del pico mensual frente al real (Países Bajos, 2024) con banda de IC 95\,\%.}
\label{fig:arima-forecast}
\end{figure}

\begin{figure}[H]
\centering
\includegraphics[width=0.82\linewidth]{var_forecast.png}
\caption{Pronóstico VAR del pico mensual frente al real (Países Bajos, 2024).}
\label{fig:var-forecast}
\end{figure}

\subsection*{Análisis de los resultados}

El ARIMA superó al VAR en las tres métricas (MAPE $2{,}64\,\%$ frente a $3{,}26\,\%$), lo que coincide con lo reportado en el artículo, donde ARIMA también fue el modelo con menor error. En series univariadas con estacionalidad marcada, un ARIMA bien calibrado puede superar a un VAR multivariado cuando las exógenas no aportan información adicional; el propio artículo ya lo señalaba al notar que el VAR rinde mejor solo con predictores externos fuertemente acoplados.

El diagnóstico de residuos muestra una diferencia: el ARIMA deja autocorrelación residual (Ljung-Box $p<0{,}05$), lo que indica que su estructura estacional no captura toda la dinámica y podría mejorarse con términos SARIMA adicionales o incluyendo exógenas (SARIMAX). El VAR, en cambio, presenta residuos sin autocorrelación ($p=0{,}62$), con un ajuste más adecuado de la estructura, aunque con mayor error de pronóstico.

Limitaciones de esta aplicación: los datos son sintéticos, con magnitudes de consumo poco realistas entre países; la muestra es corta (60 meses, 12 de prueba); y no se incluyeron los regresores Lasso, Ridge y XGBoost ni el modelo híbrido, lo cual queda como trabajo futuro para reproducir la comparación completa del artículo.
